\section{Conclusion}
\label{sec:conclusion}

This paper presented \softrobotagent{}, a research-oriented framework for tendon-driven soft-robot co-design. Its methodology combines mechanism-guided autonomous research, composable mathematical research, and evidence-grounded iteration. The central idea is to treat robot design as an adaptive scientific process: physical observations motivate hypotheses, hypotheses are translated into parameterized mathematical or physical investigations, and the resulting evidence changes subsequent research decisions.

The mathematical framework distinguishes mechanism-level, fixed-input, re-optimized task-capability, and closed-loop sensitivities. It connects established optimization and differentiation methods to a common research interface while retaining explicit validity conditions for nonconvexity, unilateral constraints, and active-set transitions.

The current task portfolio includes active multi-target control, passive progressive stiffness, passive multi-condition stabilization, tracking, and other loaded or orientation-constrained tasks. The next stage is to quantify the value of research-intent projection, mathematical guidance, pairwise interaction analysis, domain knowledge, and evidence-driven replanning through controlled comparisons. The broader goal is to move beyond LLM-assisted candidate generation toward an agent that can formulate, analyze, test, and revise scientific hypotheses about robot design.
